%% chapters/00-abstrak-id.tex
%% Abstrak Bahasa Indonesia (~500-600 kata)
%% Aturan: .cursor/rules/disertation/12-abstrak-front-matter.mdc

\begin{abstractid}

Perawatan prediktif (\emph{predictive maintenance}) berbasis kondisi
merupakan strategi kritis dalam transformasi industri manufaktur modern
menuju era Industri 4.0.
Kemampuan memprediksi sisa umur pakai (\emph{remaining useful life},
RUL) komponen kritis secara akurat dan tepat waktu memungkinkan
intervensi perawatan yang efisien, meminimalkan \emph{downtime} tidak
terencana, dan meningkatkan keselamatan operasional.
\emph{Rolling element bearings} merupakan komponen yang paling rentan terhadap
kegagalan pada mesin rotasi industri, dengan kontribusi sekitar
40--50\,\% dari total kegagalan mesin.

Disertasi ini mengusulkan arsitektur hibrida \emph{Mamba-xLSTM-Net}
untuk prediksi RUL \emph{rolling element bearings}, yang mengintegrasikan
\emph{selective state-space model} (Mamba-3) dengan \emph{matrix-memory
Long Short-Term Memory} (mLSTM).
Mekanisme fusi dirancang berdasarkan sifat non-stasioner sinyal getaran
\emph{bearing}: Mamba menangani dependensi temporal jangka panjang pada fase
degradasi gradual, sementara mLSTM mempertahankan presisi rekuren pada
perubahan mendadak menjelang kegagalan.

Di samping aspek kinerja prediksi, disertasi ini juga mengusulkan modul
interpretabilitas berbasis \emph{Top-$k$ Sparse Autoencoder} (SAE) yang
bekerja secara \emph{post-hoc} terhadap representasi laten
\emph{Mamba-xLSTM-Net}.
Prosedur kuantitatif dikembangkan untuk memetakan fitur-fitur SAE ke
frekuensi karakteristik fisik \emph{bearing} --- \emph{Ball Pass Frequency
Outer Race} (BPFO), \emph{Ball Pass Frequency Inner Race} (BPFI),
\emph{Ball Spin Frequency} (BSF), dan \emph{Fundamental Train Frequency}
(FTF) --- melalui korelasi Pearson terhadap amplitudo \emph{Hilbert
envelope spectrum} sinyal getaran mentah.

Validasi empiris dilakukan pada dua dataset publik \emph{run-to-failure},
yaitu PHM2012 (FEMTO-PRONOSTIA) yang memuat 17 \emph{bearings} NSK~6804 dalam
tiga kondisi operasi, dan XJTU-SY yang memuat 10 \emph{bearings} LDK~UER204
(kondisi 1 dan 2).
Protokol multi-\emph{seed} dengan tiga inisialisasi berbeda (seed 42, 43,
44) diterapkan untuk memastikan rigor statistik.
Kinerja dievaluasi menggunakan empat metrik sekaligus: RMSE, MAE,
koefisien determinasi R\textsuperscript{2}, dan PHM Score asimetris.

Pada dataset PHM2012, \emph{Mamba-xLSTM-Net} mencapai RMSE
$0{,}242 \pm 0{,}020$ dan memperlihatkan stabilitas lintas \emph{seed}
tertinggi di antara tiga model yang dibandingkan.
Pada XJTU-SY, \emph{Mamba-xLSTM-Net} mencapai PHM Score tertinggi
($0{,}907$), metrik yang paling relevan dengan praktik pemeliharaan
prediktif karena memperhitungkan asimetri konsekuensi kesalahan.
Kurva pelatihan yang terus membaik hingga \emph{epoch} 71/75 menunjukkan
bahwa arsitektur ini belum mencapai kapasitas penuh pada anggaran yang
digunakan.

Pada Pilar~2, SAE dengan 1024 dimensi laten dan $k = 51$ fitur aktif
($\approx 5\%$ sparsity) berhasil merekonstruksi representasi laten
dengan \emph{loss} MSE di bawah $10^{-3}$, mengindikasikan bahwa
representasi laten \emph{Mamba-xLSTM-Net} secara alami bersifat
\emph{sparse}.
Pemetaan BPFx menunjukkan bahwa 2{,}0--2{,}3\% fitur SAE berkorelasi
moderat ($|r| \geq 0{,}30$) dengan amplitudo \emph{envelope spectrum}
di sekitar BPFO dan BPFI, yang merupakan satu-satunya frekuensi
karakteristik dengan \emph{hit-rate} $> 0$.
Korelasi moderat tertinggi mencapai $r = 0{,}507$ (fitur SAE
f474~$\leftrightarrow$~BPFI pada PHM2012) dan $r = 0{,}501$
(f959~$\leftrightarrow$~BPFO pada XJTU-SY), keduanya signifikan secara
statistik ($n = 300$, $p \ll 0{,}05$).
Hasil ini merupakan bukti empiris pertama bahwa model \emph{deep
learning} yang dilatih untuk RUL mempelajari representasi laten yang
berkorespondensi dengan teori klasik diagnosis getaran \emph{bearing}.

Penelitian ini memberikan tiga kontribusi utama: (1)~arsitektur
hibrida \emph{Mamba-xLSTM-Net} sebagai pendekatan baru yang menggabungkan
SSM selektif dengan \emph{matrix-memory recurrent} untuk prognostik
\emph{bearing}; (2)~prosedur pemetaan SAE ke frekuensi karakteristik fisik
sebagai metodologi interpretabilitas baru untuk domain perawatan
prediktif; dan (3)~bukti empiris korespondensi antara representasi laten
\emph{deep learning} dan teori fisika klasik pada dua dataset
\emph{benchmark} terkemuka.

\end{abstractid}
